\textbf{Datos y mercado.} Se asume que los precios de cierre y dividendos de los CSV representan adecuadamente la informaci\'on disponible para cada activo, por lo que el retorno total incorpora reinversi\'on impl\'icita de dividendos. La matriz de retornos se construye con \texttt{join="inner"}, suponiendo que eliminar fechas no comunes no introduce un sesgo relevante. La metadata de capitalizaci\'on de mercado se considera suficiente para construir el benchmark y los pesos de equilibrio de Black-Litterman, aunque puede existir sesgo de supervivencia. Las carteras son \textit{long-only}, cumplen $w_i \geq 0$ y $\sum w_i = 1$, no usan apalancamiento y permiten divisibilidad perfecta de acciones. La optimizaci\'on omite costos de transacci\'on, impuestos, \textit{spreads} bid-ask, restricciones de liquidez, impacto de mercado y \textit{slippage}; en Monte Carlo, estos costos se aproximan con una comisi\'on de $k = 1\%$ sobre transacciones. La tasa libre de riesgo se mantiene constante en $R_f = 2\%$ anual.

\textbf{Supuestos estad\'isticos.} Los retornos hist\'oricos se usan como informaci\'on para estimar retornos y riesgos futuros, bajo estabilidad temporal parcial de $\mu$ y $\Sigma$. Para estabilizar la covarianza muestral se aplica \textit{shrinkage} de 20\%: $\Sigma_{\text{shrink}} = (1-\lambda)\Sigma + \lambda \cdot \text{diag}(\Sigma)$, con $\lambda = 0.20$. La anualizaci\'on considera 252 d\'ias h\'abiles y escala la volatilidad por $\sqrt{252}$. En la simulaci\'on Monte Carlo, los retornos semanales se aproximan mediante $\mathcal{N}(\mu_{\text{sem}},\, \sigma_{\text{sem}}^2)$, aun cuando el an\'alisis distribucional muestra colas pesadas y curtosis elevada en los retornos diarios.

\textbf{Optimizaci\'on y rebalanceo.} Los problemas formulados en CVXPY se tratan como convexos y adecuados para representar la decisi\'on de inversi\'on. Si una cota de volatilidad vuelve infactible un perfil, se relaja solo esa restricci\'on. El portafolio de tangencia se aproxima eligiendo el mayor Sharpe dentro de la frontera eficiente discretizada, usando 30 puntos en Markowitz y 12 en Black-Litterman. En Black-Litterman, las \textit{views} no requieren cubrir todos los activos: aquellos sin \textit{view} expl\'icita conservan el prior de equilibrio $\pi$ y reciben ajustes indirectos por correlaci\'on a trav\'es de $\Sigma$. Las variantes por capitalizaci\'on usan subconjuntos de 20 a 40 activos, mientras que momentum general y desempleo consideran los 601 activos. El rebalanceo se realiza cada 126 d\'ias h\'abiles; entre rebalanceos los pesos permanecen constantes y, en cada fecha, la ventana de entrenamiento incorpora solo retornos ya observados.

\textbf{Experimento pandemia.} La comparaci\'on entre escenarios ``con pandemia'' y ``sin pandemia'' mantiene fijo el per\'iodo futuro fuera de muestra. Por ello, las diferencias se interpretan como efecto de la informaci\'on utilizada para calibrar los modelos, no como cambios en el mercado futuro. La base sin pandemia excluye las observaciones de 2020--2022, mientras que la base con pandemia las conserva.

\textbf{Simulaci\'on Monte Carlo (P4).} Cada cliente inicia con $C_0 = 1\,000$ USD. La comisi\'on es $k = 1\%$, aplicada sobre $C_0$ al inicio y sobre el 5\% de la riqueza actual en cada recomendaci\'on aceptada, como fracci\'on de rotaci\'on asumida. La p\'erdida se mide respecto de $C_0$ y no del m\'aximo hist\'orico de riqueza, de modo que una ca\'ida posterior a una ganancia no activa $P_1$ mientras la riqueza siga sobre el capital inicial. La sensibilidad log\'istica $s = 20$ escala las diferencias decimales de $P_1$ y $P_2$; para el perfil muy conservador, con tolerancia 0\%, cualquier p\'erdida genera una probabilidad de retiro cercana a 1. Los escenarios proyectados se definen como $R_{\text{base}} \pm 0.5\,\sigma$, formando un abanico sim\'etrico; aunque los retornos reales pueden presentar asimetr\'ia negativa, esta aproximaci\'on es suficiente para el objetivo acad\'emico. Se ejecutan 5\,000 simulaciones por combinaci\'on, cantidad considerada adecuada para estabilizar riqueza terminal, tasa de retiro y utilidad de la empresa.

\textbf{Sensibilidad.} El impacto de los supuestos se eval\'ua mediante tres ejercicios: comparaci\'on del entrenamiento con y sin pandemia, midiendo el delta de Sharpe por portafolio y t\'ecnica; validaci\'on con 70 splits aleatorios y 18 ventanas m\'oviles robustas, con horizontes desde h2\_f7 hasta h7\_f2; y proyecciones Monte Carlo desfavorable, neutra y favorable, modificando el retorno esperado en $\pm 0.5\sigma$ para observar cambios en riqueza terminal, tasa de retiro y utilidad de la empresa.
